% --- Archon physics formalization source begin ---
% archon:physics
% archon:covers PhyXMiniProblems/problem_phyx_mini_0249.lean
% archon:source-report reports/phyx_mini/problem_phyx_mini_0249.source.json

\chapter{Physics problem phyx\_mini\_0249}
\label{ch:PhyXMiniProblems_problem_phyx_mini_0249}

\paragraph{Problem source.}
\#\# Physical scenario

Light from a helium-neon laser (\$\textbackslash{}lambda = 633\textbackslash{};nm\$) passes through a narrow slit and is seen on a screen \$2.0\textbackslash{};m\$ behind the slit. The first minimum in the diffraction pattern is \$1.2\textbackslash{};cm\$ from the central maximum.

\#\# Question

How wide is the slit?

\#\# Answer choices

- A: A: 0.11mm

- B: B: 0.13mm

- C: C: 0.10mm

- D: D: 0.12mm

\#\# Figure caption (auxiliary; use the image as primary evidence)

The image depicts a diagram of a single-slit diffraction experiment. 

- On the left, there is a vertical bar labeled "Single slit," with a width indicated as \textbackslash{}( a \textbackslash{}).
- A wave pattern representing light intensity is shown spreading out from the slit toward the right where it meets a "Screen."
- The central maximum is labeled; it is the brightest part of the pattern on the screen.
- Dashed lines show angles \textbackslash{}(\textbackslash{}theta\_1\textbackslash{}) extending from the center of the slit to the screen.
- The light intensity is depicted as a curve from the slit to the screen.
- Peaks and troughs on the graph are marked as \textbackslash{}( p = 1 \textbackslash{}) and \textbackslash{}( p = 2 \textbackslash{}) away from the central maximum.
- The distance between points marked \textbackslash{}( p = 1 \textbackslash{}) and \textbackslash{}( p = 1 \textbackslash{}) on either side of the central maximum is labeled as "Width \textbackslash{}( w \textbackslash{})."
- A note states that the distance \textbackslash{}( L \textbackslash{}gg a \textbackslash{}), indicating that the screen is far from the slit.

\#\# Dataset metadata

Wave Optics; Physical Model Grounding Reasoning, Spatial Relation Reasoning

\paragraph{Recorded answer/context.}
A: A: 0.11mm

\paragraph{Figure/image path.}
/root/proposal\_for\_physic/science-mango/phyx\_mini\_run/phyx\_data/test\_image/249.png

\paragraph{Formalization target.}
create a compiling Lean file with sorry bodies at `PhyXMiniProblems/problem\_phyx\_mini\_0249.lean`. The Lean declarations must preserve the physical quantities, dimensions or dimensional roles, figure labels, governing-law hypotheses, and final relation expressed by this problem.
Use Mathlib/Physlib names found through LeanExplore where available. If a physics API is missing, introduce faithful local abstractions rather than scalar placeholder aliases.

\begin{theorem}[Physics formalization target]
\label{thm:physics:phyx_mini_0249:target}
\lean{PhyXMiniProblems.ProblemPhyXMini0249.problem_phyx_mini_0249}
\uses{def:physics:phyx-mini-0249:phyxminiproblems-problemphyxmini0249-lengthinmillimeters, def:physics:phyx-mini-0249:phyxminiproblems-problemphyxmini0249-singleslitdiffractionsetup, def:physics:phyx-mini-0249:phyxminiproblems-problemphyxmini0249-matchesproblemandfigurereadouts, def:physics:phyx-mini-0249:phyxminiproblems-problemphyxmini0249-hasphysicalparameters, def:physics:phyx-mini-0249:phyxminiproblems-problemphyxmini0249-smallpositivescale, def:physics:phyx-mini-0249:phyxminiproblems-problemphyxmini0249-connectsexperimenttoscaledfamily, def:physics:phyx-mini-0249:phyxminiproblems-problemphyxmini0249-satisfiesexactscreenraygeometry, def:physics:phyx-mini-0249:phyxminiproblems-problemphyxmini0249-satisfiesfraunhoferparaxialasymptotics, def:physics:phyx-mini-0249:phyxminiproblems-problemphyxmini0249-widthestimateprofileinmillimeters, def:physics:phyx-mini-0249:phyxminiproblems-problemphyxmini0249-inferredslitwidthinmillimeters, def:physics:phyx-mini-0249:phyxminiproblems-problemphyxmini0249-recordeddatasetanswer, def:physics:phyx-mini-0249:phyxminiproblems-problemphyxmini0249-isuniquematchingdisplayedwidth}
The assigned autoformalize agent should translate this physics problem into Lean declarations in the covered file, with theorem and lemma proof bodies written as `by sorry`.
\end{theorem}
\begin{proof}
This is an autoformalization task, not a proof task. Produce statements that can later be proved by the physics prover without weakening the physical model.
\end{proof}

% --- Archon declaration topology begin ---
\section{Declaration topology}
\begin{definition}[Optical Length]
  \label{def:physics:phyx-mini-0249:phyxminiproblems-problemphyxmini0249-opticallength}
  \lean{PhyXMiniProblems.ProblemPhyXMini0249.OpticalLength}
  A signed, unit-independent physical quantity carrying length dimension
\end{definition}
\begin{proof}
  This is the declared model definition of Optical Length.
\end{proof}
\begin{definition}[length Readout]
  \label{def:physics:phyx-mini-0249:phyxminiproblems-problemphyxmini0249-lengthreadout}
  \lean{PhyXMiniProblems.ProblemPhyXMini0249.lengthReadout}
  \uses{def:physics:phyx-mini-0249:phyxminiproblems-problemphyxmini0249-opticallength}
  Read a physical length as a real number in a selected length unit
\end{definition}
\begin{proof}
  This is the declared model definition of length Readout.
\end{proof}
\begin{definition}[length In Meters]
  \label{def:physics:phyx-mini-0249:phyxminiproblems-problemphyxmini0249-lengthinmeters}
  \lean{PhyXMiniProblems.ProblemPhyXMini0249.lengthInMeters}
  \uses{def:physics:phyx-mini-0249:phyxminiproblems-problemphyxmini0249-opticallength, def:physics:phyx-mini-0249:phyxminiproblems-problemphyxmini0249-lengthreadout}
  Metre readout of a physical length
\end{definition}
\begin{proof}
  This is the declared model definition of length In Meters.
\end{proof}
\begin{definition}[length In Centimeters]
  \label{def:physics:phyx-mini-0249:phyxminiproblems-problemphyxmini0249-lengthincentimeters}
  \lean{PhyXMiniProblems.ProblemPhyXMini0249.lengthInCentimeters}
  \uses{def:physics:phyx-mini-0249:phyxminiproblems-problemphyxmini0249-opticallength, def:physics:phyx-mini-0249:phyxminiproblems-problemphyxmini0249-lengthreadout}
  Centimetre readout of a physical length
\end{definition}
\begin{proof}
  This is the declared model definition of length In Centimeters.
\end{proof}
\begin{definition}[length In Millimeters]
  \label{def:physics:phyx-mini-0249:phyxminiproblems-problemphyxmini0249-lengthinmillimeters}
  \lean{PhyXMiniProblems.ProblemPhyXMini0249.lengthInMillimeters}
  \uses{def:physics:phyx-mini-0249:phyxminiproblems-problemphyxmini0249-opticallength, def:physics:phyx-mini-0249:phyxminiproblems-problemphyxmini0249-lengthreadout}
  Millimetre readout of a physical length
\end{definition}
\begin{proof}
  This is the declared model definition of length In Millimeters.
\end{proof}
\begin{definition}[length In Nanometers]
  \label{def:physics:phyx-mini-0249:phyxminiproblems-problemphyxmini0249-lengthinnanometers}
  \lean{PhyXMiniProblems.ProblemPhyXMini0249.lengthInNanometers}
  \uses{def:physics:phyx-mini-0249:phyxminiproblems-problemphyxmini0249-opticallength, def:physics:phyx-mini-0249:phyxminiproblems-problemphyxmini0249-lengthreadout}
  Nanometre readout of a physical length
\end{definition}
\begin{proof}
  This is the declared model definition of length In Nanometers.
\end{proof}
\begin{definition}[Laser Kind]
  \label{def:physics:phyx-mini-0249:phyxminiproblems-problemphyxmini0249-laserkind}
  \lean{PhyXMiniProblems.ProblemPhyXMini0249.LaserKind}
  The laser source named in the problem
\end{definition}
\begin{proof}
  This is the declared model definition of Laser Kind.
\end{proof}
\begin{definition}[Slit Geometry]
  \label{def:physics:phyx-mini-0249:phyxminiproblems-problemphyxmini0249-slitgeometry}
  \lean{PhyXMiniProblems.ProblemPhyXMini0249.SlitGeometry}
  The aperture geometry shown on the left of the supplied diagram
\end{definition}
\begin{proof}
  This is the declared model definition of Slit Geometry.
\end{proof}
\begin{definition}[Optical Regime]
  \label{def:physics:phyx-mini-0249:phyxminiproblems-problemphyxmini0249-opticalregime}
  \lean{PhyXMiniProblems.ProblemPhyXMini0249.OpticalRegime}
  The controlled limiting model used for the diffraction calculation
\end{definition}
\begin{proof}
  This is the declared model definition of Optical Regime.
\end{proof}
\begin{definition}[Figure Feature]
  \label{def:physics:phyx-mini-0249:phyxminiproblems-problemphyxmini0249-figurefeature}
  \lean{PhyXMiniProblems.ProblemPhyXMini0249.FigureFeature}
  Text labels and qualitative features visible in the supplied figure
\end{definition}
\begin{proof}
  This is the declared model definition of Figure Feature.
\end{proof}
\begin{definition}[Single Slit Diffraction Setup]
  \label{def:physics:phyx-mini-0249:phyxminiproblems-problemphyxmini0249-singleslitdiffractionsetup}
  \lean{PhyXMiniProblems.ProblemPhyXMini0249.SingleSlitDiffractionSetup}
  \uses{def:physics:phyx-mini-0249:phyxminiproblems-problemphyxmini0249-opticallength, def:physics:phyx-mini-0249:phyxminiproblems-problemphyxmini0249-laserkind, def:physics:phyx-mini-0249:phyxminiproblems-problemphyxmini0249-slitgeometry, def:physics:phyx-mini-0249:phyxminiproblems-problemphyxmini0249-opticalregime, def:physics:phyx-mini-0249:phyxminiproblems-problemphyxmini0249-figurefeature}
  This structure records the Single Slit Diffraction Setup component of the problem-specific physical model
\end{definition}
\begin{proof}
  This is the declared model definition of Single Slit Diffraction Setup.
\end{proof}
\begin{definition}[Matches Problem And Figure Readouts]
  \label{def:physics:phyx-mini-0249:phyxminiproblems-problemphyxmini0249-matchesproblemandfigurereadouts}
  \lean{PhyXMiniProblems.ProblemPhyXMini0249.MatchesProblemAndFigureReadouts}
  \uses{def:physics:phyx-mini-0249:phyxminiproblems-problemphyxmini0249-lengthreadout, def:physics:phyx-mini-0249:phyxminiproblems-problemphyxmini0249-lengthinmeters, def:physics:phyx-mini-0249:phyxminiproblems-problemphyxmini0249-lengthincentimeters, def:physics:phyx-mini-0249:phyxminiproblems-problemphyxmini0249-lengthinnanometers, def:physics:phyx-mini-0249:phyxminiproblems-problemphyxmini0249-singleslitdiffractionsetup}
  Source values and primary-figure readouts. The full central-maximum width is the distance between the two symmetric first minima, hence w = 2 y₁. This predicate contains no value of the unknown slit width or of a displayed answer
\end{definition}
\begin{proof}
  This is the declared model definition of Matches Problem And Figure Readouts.
\end{proof}
\begin{definition}[Has Physical Parameters]
  \label{def:physics:phyx-mini-0249:phyxminiproblems-problemphyxmini0249-hasphysicalparameters}
  \lean{PhyXMiniProblems.ProblemPhyXMini0249.HasPhysicalParameters}
  \uses{def:physics:phyx-mini-0249:phyxminiproblems-problemphyxmini0249-lengthinmeters, def:physics:phyx-mini-0249:phyxminiproblems-problemphyxmini0249-singleslitdiffractionsetup}
  Positive, nondegenerate parameters of the source experiment
\end{definition}
\begin{proof}
  This is the declared model definition of Has Physical Parameters.
\end{proof}
\begin{definition}[Is Positive Acute Angle]
  \label{def:physics:phyx-mini-0249:phyxminiproblems-problemphyxmini0249-ispositiveacuteangle}
  \lean{PhyXMiniProblems.ProblemPhyXMini0249.IsPositiveAcuteAngle}
  An angle lies on the positive acute branch when it has a representative strictly between 0 and π/2
\end{definition}
\begin{proof}
  This is the declared model definition of Is Positive Acute Angle.
\end{proof}
\begin{definition}[small Positive Scale]
  \label{def:physics:phyx-mini-0249:phyxminiproblems-problemphyxmini0249-smallpositivescale}
  \lean{PhyXMiniProblems.ProblemPhyXMini0249.smallPositiveScale}
  Approach zero through positive values of the approximation parameter
\end{definition}
\begin{proof}
  This is the declared model definition of small Positive Scale.
\end{proof}
\begin{definition}[Connects Experiment To Scaled Family]
  \label{def:physics:phyx-mini-0249:phyxminiproblems-problemphyxmini0249-connectsexperimenttoscaledfamily}
  \lean{PhyXMiniProblems.ProblemPhyXMini0249.ConnectsExperimentToScaledFamily}
  \uses{def:physics:phyx-mini-0249:phyxminiproblems-problemphyxmini0249-lengthreadout, def:physics:phyx-mini-0249:phyxminiproblems-problemphyxmini0249-lengthinmillimeters, def:physics:phyx-mini-0249:phyxminiproblems-problemphyxmini0249-singleslitdiffractionsetup}
  This interface connects the actual experiment to a scaled family. The family keeps the physical slit width fixed and uses λ(ε) = ε a; thus ε has the physical role of the dimensionless ratio λ / a. It also identifies the measured first-minimum ratio at the selected experimental scale. These are parameterization and observation relations, not the numerical answer
\end{definition}
\begin{proof}
  This is the declared model definition of Connects Experiment To Scaled Family.
\end{proof}
\begin{definition}[Satisfies Exact Screen Ray Geometry]
  \label{def:physics:phyx-mini-0249:phyxminiproblems-problemphyxmini0249-satisfiesexactscreenraygeometry}
  \lean{PhyXMiniProblems.ProblemPhyXMini0249.SatisfiesExactScreenRayGeometry}
  \uses{def:physics:phyx-mini-0249:phyxminiproblems-problemphyxmini0249-singleslitdiffractionsetup, def:physics:phyx-mini-0249:phyxminiproblems-problemphyxmini0249-ispositiveacuteangle}
  Exact planar-screen ray geometry uses yₚ / L = tan θₚ. This equality is geometrical and is not the paraxial approximation tan θ ≈ sin θ
\end{definition}
\begin{proof}
  This is the declared model definition of Satisfies Exact Screen Ray Geometry.
\end{proof}
\begin{definition}[Satisfies Fraunhofer Paraxial Asymptotics]
  \label{def:physics:phyx-mini-0249:phyxminiproblems-problemphyxmini0249-satisfiesfraunhoferparaxialasymptotics}
  \lean{PhyXMiniProblems.ProblemPhyXMini0249.SatisfiesFraunhoferParaxialAsymptotics}
  \uses{def:physics:phyx-mini-0249:phyxminiproblems-problemphyxmini0249-lengthinmillimeters, def:physics:phyx-mini-0249:phyxminiproblems-problemphyxmini0249-singleslitdiffractionsetup, def:physics:phyx-mini-0249:phyxminiproblems-problemphyxmini0249-smallpositivescale}
  A controlled asymptotic form of the two approximations used by the textbook calculation. The single-slit residual a sin θₚ(ε) - p λ(ε) is little-o of λ(ε). The paraxial projection residual yₚ(ε)/L - sin θₚ(ε) is little-o of sin θₚ(ε). Consequently neither approximation is asserted as a global exact identity at the finite source experiment. The exact screen relation remains the tangent law above
\end{definition}
\begin{proof}
  This is the declared model definition of Satisfies Fraunhofer Paraxial Asymptotics.
\end{proof}
\begin{definition}[width Estimate Profile In Millimeters]
  \label{def:physics:phyx-mini-0249:phyxminiproblems-problemphyxmini0249-widthestimateprofileinmillimeters}
  \lean{PhyXMiniProblems.ProblemPhyXMini0249.widthEstimateProfileInMillimeters}
  \uses{def:physics:phyx-mini-0249:phyxminiproblems-problemphyxmini0249-singleslitdiffractionsetup}
  The order-one Fraunhofer-paraxial width estimator along the scaled family. Both numerator and result are millimetre scalar readouts, while the denominator is the dimensionless ratio y₁ / L
\end{definition}
\begin{proof}
  This is the declared model definition of width Estimate Profile In Millimeters.
\end{proof}
\begin{definition}[inferred Slit Width In Millimeters]
  \label{def:physics:phyx-mini-0249:phyxminiproblems-problemphyxmini0249-inferredslitwidthinmillimeters}
  \lean{PhyXMiniProblems.ProblemPhyXMini0249.inferredSlitWidthInMillimeters}
  \uses{def:physics:phyx-mini-0249:phyxminiproblems-problemphyxmini0249-lengthinmillimeters, def:physics:phyx-mini-0249:phyxminiproblems-problemphyxmini0249-singleslitdiffractionsetup}
  The leading-order estimator evaluated directly from the three source measurements: a\_est = λ L / y₁. This is a general measurement formula, not a definition of the requested numerical answer or of the physical slit width
\end{definition}
\begin{proof}
  This is the declared model definition of inferred Slit Width In Millimeters.
\end{proof}
\begin{definition}[Answer Choice]
  \label{def:physics:phyx-mini-0249:phyxminiproblems-problemphyxmini0249-answerchoice}
  \lean{PhyXMiniProblems.ProblemPhyXMini0249.AnswerChoice}
  Labels of the four slit-width choices printed with the dataset item
\end{definition}
\begin{proof}
  This is the declared model definition of Answer Choice.
\end{proof}
\begin{definition}[displayed Width In Millimeters]
  \label{def:physics:phyx-mini-0249:phyxminiproblems-problemphyxmini0249-displayedwidthinmillimeters}
  \lean{PhyXMiniProblems.ProblemPhyXMini0249.displayedWidthInMillimeters}
  \uses{def:physics:phyx-mini-0249:phyxminiproblems-problemphyxmini0249-answerchoice}
  Millimetre value printed beside each answer label
\end{definition}
\begin{proof}
  This is the declared model definition of displayed Width In Millimeters.
\end{proof}
\begin{definition}[recorded Dataset Answer]
  \label{def:physics:phyx-mini-0249:phyxminiproblems-problemphyxmini0249-recordeddatasetanswer}
  \lean{PhyXMiniProblems.ProblemPhyXMini0249.recordedDatasetAnswer}
  \uses{def:physics:phyx-mini-0249:phyxminiproblems-problemphyxmini0249-answerchoice}
  The answer label recorded in the supplied dataset
\end{definition}
\begin{proof}
  This is the declared model definition of recorded Dataset Answer.
\end{proof}
\begin{definition}[rounded To Nearest Hundredth]
  \label{def:physics:phyx-mini-0249:phyxminiproblems-problemphyxmini0249-roundedtonearesthundredth}
  \lean{PhyXMiniProblems.ProblemPhyXMini0249.roundedToNearestHundredth}
  Round a millimetre readout to the nearest hundredth of a millimetre
\end{definition}
\begin{proof}
  This is the declared model definition of rounded To Nearest Hundredth.
\end{proof}
\begin{definition}[Matches Displayed Width]
  \label{def:physics:phyx-mini-0249:phyxminiproblems-problemphyxmini0249-matchesdisplayedwidth}
  \lean{PhyXMiniProblems.ProblemPhyXMini0249.MatchesDisplayedWidth}
  \uses{def:physics:phyx-mini-0249:phyxminiproblems-problemphyxmini0249-answerchoice, def:physics:phyx-mini-0249:phyxminiproblems-problemphyxmini0249-displayedwidthinmillimeters, def:physics:phyx-mini-0249:phyxminiproblems-problemphyxmini0249-roundedtonearesthundredth}
  A displayed choice agrees with an inferred width at the shown precision
\end{definition}
\begin{proof}
  This is the declared model definition of Matches Displayed Width.
\end{proof}
\begin{definition}[Is Unique Matching Displayed Width]
  \label{def:physics:phyx-mini-0249:phyxminiproblems-problemphyxmini0249-isuniquematchingdisplayedwidth}
  \lean{PhyXMiniProblems.ProblemPhyXMini0249.IsUniqueMatchingDisplayedWidth}
  \uses{def:physics:phyx-mini-0249:phyxminiproblems-problemphyxmini0249-answerchoice, def:physics:phyx-mini-0249:phyxminiproblems-problemphyxmini0249-matchesdisplayedwidth}
  The selected label is the unique displayed choice matching the estimate
\end{definition}
\begin{proof}
  This is the declared model definition of Is Unique Matching Displayed Width.
\end{proof}
% --- Archon declaration topology end ---
% --- Archon physics formalization source end ---
